\section{Introduction}
\subsection{Chapter subdivision}
\paragraph{The anti-neutron in physics:} the Chapter 2 introduces the anti-neutron in the high-energy physics (HEP) context. Its primary interactions -such as annihilation and hadronic shower production- are then exploited, discussing some hadronic models very quickly. Eventually, the MANTRA (Measuring Anti-Neutron: Tagging and Reconstruction Algorithm) project is then described as a method to measure anti-neutron energies up to 10 GeV.

\paragraph{The Belle II experiment:} the Chapter 3 provides an overview of the SuperKEKB accelerator and Belle II detector. In the context of $\bar{n}$ studies, the Electromagnetic Calorimeter (ECL) plays a key role in hadron detection. A detailed description of the ECL is thus provided, covering its inorganic crystals, reconstruction algorithms, and performance. 
$\\$
Then the Belle II Analysis Software Framework (basf2) is introduced, providing an overview of its role in the simulation and reconstruction processes. Eventually the Topology Analysis (TopoAna) independent  software is described.

\paragraph{Study of the signal sample channel:} the Chapter 4 is devoted to the study of the kinematic variables of anti-neutrons $\bar{n}$, including momentum, angular distributions and related observables. The analysis follows a semi-exclusive approach over the selected channel $e^+ e^- \rightarrow p \pi^-  \bar{n}  (\gamma_{\rm ISR})$, studying both ISR and NO ISR cases, where the recoil system is made of  $p  \pi^-  (\gamma_{ISR})$. Therefore the recoil system variables are compared to those of the anti-neutrons candidates, studying the response of a Kinematic Fit as well. 
$\\$
Cluster split-off and pre-showering effects are investigated using ancestor meta-variables. Eventually several neutral cluster variables associated to anti-neutron candidates are then described and studied for both ISR and NO ISR cases, including the analysis of their ancestor contributions.



\subsection{The state of subatomic physics today}
The Standard Model is one of the greatest achievements of modern science, especially in the field of high-energy particle physics. Its reliability is the result of a long process that began in the 20th century and received its definitive confirmation in 2012 with the discovery of the Higgs boson. 
\\
However, there are some topics which remain unsolved, like the non-trivial neutrino mass, the presence of dark matter on cosmological scales or the study of heavy hadrons. The Standard Model still represents an excellent phenomenological description of subatomic processes, at the $\sim$1 TeV energy scale.
\\
In this framework, the \textbf{Belle II} experiment, located at the SuperKEKB facility, plays a fundamental role in exploring the aforementioned topics. These are scenarios that go beyond the standard physics, thus entering the so-called \emph{new physics}.
\\
In parallel with experiments that maximize the energy, the aim of Belle II is to explore the intensity frontier. Its main goal is not to increase the center-of-mass energy, but rather to explore scenarios of high instantaneous luminosity up to $8 \times10^{35} cm^{-2} s^{-1}$, around (and not limited to) the $\Upsilon(4S)$ resonance. This approach allows the study of extremely rare phenomena and the identification of potential discrepancies with the predictions of the Standard Model.
